\chapter*{Preface}
\addcontentsline{toc}{chapter}{Preface}

Algol 68 was designed in the late 1960s, and many of its ideas are now
everyday features of programming languages: every construct is an
expression, types can be combined freely, unions are tagged and checked
at run time, operators can be defined and overloaded, and the program can
run parts of itself concurrently. What kept many programmers away from the
language was less its ideas than the way they were described. The defining
document, the Revised Report on the Algorithmic Language Algol 68, has a
vocabulary of its own: a type is a \emph{mode}, a variable is a
\emph{name}, an expression is a \emph{unit}, an implicit conversion is a
\emph{coercion}. Most textbooks of the time followed the Report.

This book describes the language of the Revised Report in the words a
programmer uses today. It assumes that the reader can program in some
current language; it does not assume any particular one. Where a concept
is hard, such as references and their lifetimes, it is explained as
carefully as it needs to be, but it is not made harder by an unfamiliar
name.

The Report's own terms are not hidden, though. Error messages of the
compilers use them, and so do the Report and the older literature. Where a
concept is introduced, its name in the Report is given once, set in a
sans-serif font like \rr{this}; the index lists these names, and
Appendix~\ref{app-glossary} is a glossary of all of them.

\section*{The implementations}

The programs in this book are written for two implementations:
\begin{itemize}
\item \emph{ga68}, the Algol 68 front end of the GNU Compiler
  Collection~\cite{ga68}, which compiles to native code, and
\item \emph{a68j}~\cite{a68j}, which compiles the same language to the
  Java virtual machine.
\end{itemize}
Both accept the same language, apart from a few features that only one
of them provides, and give the same results. Every example program in
this book was run when the book was built, and the output shown after a
program is the output of that run. Most examples were run with both
ga68 and a68j; the few that use a feature one of them lacks
(semaphores, formats, \code|bytes|) were run with Algol 68
Genie and, where it has the feature, with a68j, and the text says so.
The examples of Appendix~\ref{app-nest}, which calls C, were run with
ga68 only, and those of Appendix~\ref{app-partial}, which describes an
extension that ga68 does not implement, with Algol 68 Genie and a68j.
Where the two implementations, or Algol 68 Genie~\cite{genie}, a widely
used interpreter, differ from the Report or from each other, a note says so;
Appendix~\ref{app-impl} collects these notes.

The versions used for this book are ga68 from the development
version of GCC 17.0.0 of 4 November 2026 (commit 7ced41f4d438), a68j \bookversion{} on OpenJDK 25, and Algol 68
Genie 3.13.4.

\section*{How the book is organized}

Chapters~\ref{ch-intro} to~\ref{ch-procs} cover the core of the language: values and variables,
expressions and control flow, arrays and records, functions and
operators. They are enough to write substantial programs. Chapters~\ref{ch-moretypes} to~\ref{ch-transput}
cover the rest: more built-in types, the advanced features of the type
system, concurrency, and input and output. Each chapter ends with
exercises, whose answers are at the end of the book, and with programming
problems, which have no single answer.

The model for this book was Andrew McGettrick's textbook \emph{ALGOL 68:
A First and Second Course}~\cite{mcgettrick}, written for students in the
years after the Revised Report appeared. The choice of topics and their
order follow it closely, and many of the exercises and programming
problems are modeled on it. The Revised
Report~\cite{rr} is the authority on every question of the language; the
bibliography at the end of the book lists it together with the other
works a reader may want to consult.
